% ══════════════════════════════════════════════════════════════════════════════
%  EXPERIMENT BOOKLET LAYER (Nov 2026)
%  From 2027 every Year 7 unit collects its experiments into one A4 experiment
%  booklet that puts the experimental techniques and skills first. Every skill
%  comes from the faculty Science Skills continuum (Years 7--8 column) and is
%  named the same way everywhere: a code (QP1, PC4, PA2 …) on a chip coloured by
%  its sub-strand. Copy this block into the next unit's experiment booklet.
%
%  Each experiment runs in the same order:
%    \experimentdiv{title}{summary}{lesson}\label{sec:N}
%    key vocabulary · learning intentions · \begin{skillsfocus}
%    \exppart*{Technique spotlight}{Blue}   (technique, worked example, We do)
%    \exppart{Plan and safety}{Teal}        (aim, prediction, variables, risks)
%    \exppart{Method}{Teal}
%    \exppart{Results}{Amber}
%    \exppart{Analysis and evaluation}{Coral}
%    \exppart{Skill check}{NavyBlue}
% ══════════════════════════════════════════════════════════════════════════════
\usetikzlibrary{arrows.meta,calc,positioning}
\usepackage{pifont}

\newenvironment{keepext}{\par\begin{tcolorbox}[blankest, breakable=false]}{\end{tcolorbox}\par}
% graph-paper box: \graphgrid{height}
\newcommand{\graphgrid}[1]{%
  \begin{tcolorbox}[enhanced, breakable=false, colback=white, colframe=MidGrey, boxrule=0.5pt, height=#1,
    underlay={\begin{tcbclipinterior}\draw[step=5mm, Line, very thin] (interior.south west) grid (interior.north east);
      \draw[step=1cm, MidGrey!70, thin] (interior.south west) grid (interior.north east);\end{tcbclipinterior}}]
  \end{tcolorbox}}

% ── Skills (faculty Science Skills continuum, Years 7–8) ─────────────────────
% \defskill{code}{name}{colour}{VC code}{what it means}
\makeatletter
\newcommand{\defskill}[5]{%
  \expandafter\gdef\csname sk@n@#1\endcsname{#2}%
  \expandafter\gdef\csname sk@c@#1\endcsname{#3}%
  \expandafter\gdef\csname sk@v@#1\endcsname{#4}%
  \expandafter\gdef\csname sk@d@#1\endcsname{#5}}
\newcommand{\skname}[1]{\csname sk@n@#1\endcsname}
\newcommand{\skcode}[1]{\csname sk@v@#1\endcsname}
\newcommand{\skdesc}[1]{\csname sk@d@#1\endcsname}
% \sk{PC4}: the chip — code and name on the sub-strand colour
\newcommand{\sk}[1]{%
  \edef\sk@col{\csname sk@c@#1\endcsname}%
  \tikz[baseline=(c.base)]\node[rounded corners=2.5pt, fill=\sk@col!16, draw=\sk@col, line width=0.6pt,
    text=\sk@col!55!black, inner xsep=4pt, inner ysep=2pt, font=\headingfont\bfseries\scriptsize]
    (c) {#1\,\textperiodcentered\,\skname{#1}};}
% \skr{PC4}: the chip flush right at the end of a paragraph; drops to its own line if it does not fit
\newcommand{\skr}[1]{\unskip\nobreak\hfill\penalty50\hskip1em\null\nobreak\hfill\sk{#1}}
% \skc{PC4}: the code alone, for tight table cells
\newcommand{\skc}[1]{%
  \edef\sk@col{\csname sk@c@#1\endcsname}%
  \tikz[baseline=(c.base)]\node[rounded corners=2.5pt, fill=\sk@col!16, draw=\sk@col, line width=0.6pt,
    text=\sk@col!55!black, inner xsep=3pt, inner ysep=2pt, font=\headingfont\bfseries\scriptsize]
    (c) {#1};}
\makeatother

% Questioning and predicting — VC2S8I01
\defskill{QP1}{Testable hypothesis}{Teal}{VC2S8I01}{A prediction an experiment can test, naming both variables and the direction of change: \emph{If the [IV] is \dots, the [DV] will \dots}; a reason is optional}
\defskill{QP2}{Formulate an aim}{Teal}{VC2S8I01}{One sentence naming both variables: \emph{To investigate the effect of changing [IV] on [DV]}}
% Planning and conducting — VC2S8I02, VC2S8I03
\defskill{PC1}{Plan and follow a method}{Blue}{VC2S8I02}{Follow, and later write, numbered steps in order: named equipment, quantities with units, no personal pronouns and a labelled diagram, so that someone else could repeat it (repeatable)}
\defskill{PC2}{Identify variables}{Blue}{VC2S8I02}{Name the independent variable (what you change), the dependent variable (what you measure) and at least two controlled variables (what you keep the same)}
\defskill{PC3}{Risks and safety}{Blue}{VC2S8I02}{Name each risk, how to prevent it and how to respond if it happens; follow the lab rules}
\defskill{PC4}{Precision}{Blue}{VC2S8I02}{Choose equipment with a fine enough scale, read it to its smallest division, and repeat measurements to see how closely they agree}
\defskill{PC5}{Ethics}{Blue}{VC2S8I02}{Treat people, living things, the environment and other people's data with respect}
\defskill{PC6}{Collect data and measure}{Blue}{VC2S8I02}{Take each measurement at the right time, with its unit, and record it straight away}
\defskill{PC7}{Use and select equipment}{Blue}{VC2S8I03}{Name each piece of equipment, use it correctly and choose the right piece for the job}
\defskill{PC8}{Use digital tools}{Blue}{VC2S8I03}{Use stopwatches, data loggers, sensors and spreadsheets to collect and record data}
% Processing, modelling and analysing — VC2S8I04, VC2S8I05
\defskill{PA1}{Tables}{GoldDark}{VC2S8I04}{Title, Headings, Units in the headings only, Decimals kept the same (THUD), with the independent variable in the first column}
\defskill{PA2}{Graphing}{GoldDark}{VC2S8I04}{The right type of graph, a title, axes labelled with units, the independent variable on the x-axis, an even scale and accurate points, then join the points or draw a line of best fit}
\defskill{PA3}{Mathematical relationships}{GoldDark}{VC2S8I04}{Process data with calculations: differences, averages, rates and formulas}
\defskill{PA4}{Creating models}{GoldDark}{VC2S8I04}{Build a physical, drawn or mathematical model that represents something real}
\defskill{PA5}{Patterns and trends}{GoldDark}{VC2S8I05}{Describe how the dependent variable changes as the independent variable changes, quoting data}
\defskill{PA6}{Correlation and causality}{GoldDark}{VC2S8I05}{Say whether two variables change together, and whether there is evidence that one causes the other}
\defskill{PA7}{Identify anomalies}{GoldDark}{VC2S8I05}{Spot a result that does not fit the pattern and suggest why it happened}
\defskill{PA8}{Interpreting models}{GoldDark}{VC2S8I05}{Use a model to explain an observation, and say what the model shows well and what it leaves out}
% Evaluating — VC2S8I06, VC2S8I07
\defskill{EV1}{Results against hypothesis}{OrangeDark}{VC2S8I06}{State the general trend and whether it supports the hypothesis (supported, never proven), with data to back the trend}
\defskill{EV2}{Compare with known concepts}{OrangeDark}{VC2S8I06}{Compare your results with accepted values and scientific ideas}
\defskill{EV3}{Error sources}{OrangeDark}{VC2S8I06}{Name what could have made results less accurate or less precise, and suggest how to improve the method}
\defskill{EV4}{Quality of secondary sources}{OrangeDark}{VC2S8I07}{Judge whether information from a website, book or video can be trusted}
\defskill{EV5}{Questions for further investigation}{OrangeDark}{VC2S8I07}{Pose a new question that follows from your results}
% Communicating — VC2S8I08
\defskill{CO1}{Scientific writing}{PurpleDark}{VC2S8I08}{Write in the third person, legibly, with correct spelling and grammar, short and to the point, using scientific terms}
\defskill{CO2}{Presentation formats}{PurpleDark}{VC2S8I08}{Choose the right format to show findings: a table, graph, scientific diagram, poster, report or talk}

% Skills focus panel at the top of each experiment:
%   \begin{skillsfocus} \focus{PC4}{what you do with it here} … \end{skillsfocus}
\newtcolorbox{skillsfocus}{
  enhanced, breakable=false, arc=2.5mm, boxrule=0.6pt,
  colback=white, colframe=NavyBlue!55,
  fonttitle=\headingfont\bfseries, coltitle=white, colbacktitle=NavyBlue,
  title={Skills focus --- the science skills this experiment builds},
  left=4mm, right=3.5mm, top=2mm, bottom=2mm}
\newcommand{\focus}[2]{\par\noindent\makebox[1.25cm][l]{\skc{#1}}\parbox[t]{\dimexpr\linewidth-1.25cm\relax}{\small\textbf{\skname{#1}}\,---\,#2}\par\vspace{2pt}}

% Learning intentions panel: \begin{lipanel} \inquiryli{…} \understandli{…} \sclist{…} \end{lipanel}
\newtcolorbox{lipanel}{
  enhanced, breakable=false, arc=2.5mm, boxrule=0.5pt,
  colback=LightBlue, colframe=Blue!40,
  fonttitle=\headingfont\bfseries, coltitle=MidBlue, title={Learning intentions},
  left=4mm, right=3.5mm, top=2mm, bottom=2mm}
\newcommand{\inquiryli}[1]{\par\noindent{\headingfont\bfseries\color{NavyBlue}Science Inquiry:}\ #1\par\vspace{2pt}}
\newcommand{\understandli}[1]{\par\noindent{\headingfont\bfseries\color{NavyBlue}Science Understanding:}\ #1\par\vspace{2pt}}
\newcommand{\sctitle}{\par\noindent{\headingfont\bfseries\color{NavyBlue}Success criteria}\par}

% Technique spotlight: \begin{technique}{title}{PC4} … \end{technique}
\newtcolorbox{technique}[2]{
  enhanced, breakable, arc=2.5mm, boxrule=0.6pt,
  colback=white, colframe=Blue!55,
  fonttitle=\headingfont\bfseries, coltitle=MidBlue, colbacktitle=LightBlue,
  title={Technique: #1\hfill\sk{#2}},
  left=4mm, right=3.5mm, top=2.5mm, bottom=2.5mm}

% Experiment divider: \experimentdiv{title}{summary}{lesson number}
% Steps the section counter, so tasks, worked examples and figures are N.M.
\newcommand{\experimentdiv}[3]{%
  \clearpage
  \refstepcounter{section}%
  \pdfbookmark[1]{Experiment \thesection{} — #1}{exp.\thesection}%
  \begin{tcolorbox}[
    enhanced, breakable=false, arc=3.5mm,
    colback=NavyBlue, colframe=NavyBlue, boxrule=0pt,
    sidebyside, sidebyside align=center,
    lefthand width=2.6cm,
    left=2mm, right=4mm, top=4mm, bottom=4mm]
    \begin{center}
      {\headingfont\bfseries\color{white!60!NavyBlue}\scriptsize EXPERIMENT}\\[-2pt]
      {\headingfont\bfseries\color{white!35!NavyBlue}\fontsize{46}{46}\selectfont \thesection}
    \end{center}
    \tcblower
    {\headingfont\bfseries\color{white}\LARGE #1}\\[3pt]
    {\itshape\color{white!75!NavyBlue}\small #2}\\[3pt]
    {\headingfont\color{Sky}\footnotesize Lesson #3}
  \end{tcolorbox}%
}

% Part bar inside an experiment. \exppart{title}{colour} starts a new page
% (new parts always start on a new page); \exppart*{…}{…} follows straight on.
\NewDocumentCommand{\exppart}{s m m}{%
  \par\IfBooleanTF{#1}{\needspace{10\baselineskip}\vspace{8pt}}{\clearpage}%
  \begin{tcolorbox}[enhanced, breakable=false, arc=2mm, boxrule=0pt,
    colback=#3!14, borderline west={4pt}{0pt}{#3},
    left=4mm, right=3mm, top=1.6mm, bottom=1.6mm]
  {\headingfont\bfseries\color{#3!60!black}\large #2}
  \end{tcolorbox}\vspace{2pt}}

% Skill check at the end of each experiment: students tick one column per row.
%   \begin{skillcheck} \scrow{PC4}{I can …} … \end{skillcheck}
\newcommand{\tickbox}{\tikz[baseline=-0.6ex]\draw[DarkGrey, line width=0.7pt, rounded corners=1pt] (0,-0.17) rectangle (0.34,0.17);}
\newenvironment{skillcheck}{%
  \par\noindent Read each statement. Tick the column that best describes you \emph{today}.\par
  \begin{ctab}\renewcommand{\arraystretch}{1.25}%
  \begin{tabular}{|W{3.7cm}|W{6.7cm}|Z{1.45cm}|Z{1.45cm}|Z{1.45cm}|}
  \hline\thead{Skill} & \thead{Success criterion} & \thead{Not yet} & \thead{Nearly} & \thead{Got it} \\ \hline}
  {\end{tabular}\end{ctab}}
\newcommand{\scrow}[2]{\rule{0pt}{1.1cm}\skc{#1}\enspace{\small\skname{#1}} & #2 & \tickbox & \tickbox & \tickbox \\ \hline}

% Write-in label under a drawing: \nameline{label}
\newcommand{\nameline}[1]{\par\noindent{\small #1}\ \leaders\hbox{\rule[-2pt]{1pt}{0.4pt}}\hfill\null\par}

% ── Equipment drawn as scientific diagrams (2D, side-on, outline only) ──────
% Each \eq... macro draws at the origin inside a tikzpicture. \eqpic[scale]{\eqBeaker}
% puts one on its own. The same shapes are used everywhere, so the booklet
% models exactly the drawings students are asked to make.
\tikzset{
  eq/.style={line width=0.9pt, draw=Ink, line cap=round, line join=round},
  eqlab/.style={draw=Ink, line width=0.4pt},
  eqtext/.style={font=\footnotesize, text=Ink, inner sep=1.5pt},
}
\newcommand{\eqpic}[2][1]{\begin{tikzpicture}[scale=#1, baseline=(current bounding box.center)]#2\end{tikzpicture}}
\newcommand{\eqBeaker}{\draw[eq] (-0.22,2.62) -- (0,2.45) -- (0,0) -- (1.8,0) -- (1.8,2.5);}
\newcommand{\eqFlask}{\draw[eq] (0.62,2.9) -- (0.65,2.8) -- (0.65,1.9) -- (0,0) -- (1.9,0) -- (1.25,1.9) -- (1.25,2.8) -- (1.28,2.9);}
\newcommand{\eqTestTube}{\draw[eq] (-0.06,3.0) -- (0,2.9) -- (0,0.3) arc[start angle=180,end angle=360,radius=0.3] -- (0.6,2.9) -- (0.66,3.0);}
\newcommand{\eqCylinder}{%
  \draw[eq] (-0.15,3.6) -- (0,3.5) -- (0,0.15) (0.6,0.15) -- (0.6,3.5);
  \draw[eq] (-0.35,0) rectangle (0.95,0.15);
  \foreach \y in {0.6,1.1,...,3.2} \draw[Ink, line width=0.5pt] (0.6,\y) -- (0.4,\y);}
\newcommand{\eqFunnel}{\draw[eq] (0,2.2) -- (0.75,1.0) -- (0.75,0) (1.85,2.2) -- (1.1,1.0) -- (1.1,0);}
\newcommand{\eqBasin}{\draw[eq] (0,0.9) .. controls (0.15,0.02) and (2.25,0.02) .. (2.4,0.9);}
\newcommand{\eqTripod}{%
  \draw[eq] (-0.15,2.0) rectangle (2.55,2.12);
  \draw[eq] (0.25,2.0) -- (0,0) (2.15,2.0) -- (2.4,0) (1.2,2.0) -- (1.2,0.15);}
\newcommand{\eqBunsen}{%
  \draw[eq] (0,0) rectangle (1.6,0.22);
  \draw[eq] (0.58,0.85) -- (0.58,2.3) (1.02,0.85) -- (1.02,2.3);
  \draw[eq] (0.5,0.22) rectangle (1.1,0.85);
  \fill[Ink] (0.72,0.42) rectangle (0.88,0.62);
  \draw[eq] (-0.65,0.06) rectangle (0,0.16);}
\newcommand{\eqThermometer}{%
  \draw[eq] (0.25,0.25) circle (0.2);
  \draw[eq] (0.15,0.43) -- (0.15,3.9) arc[start angle=180,end angle=0,radius=0.1] -- (0.35,0.43);
  \draw[Ink, line width=1.3pt] (0.25,0.45) -- (0.25,1.9);}
\newcommand{\eqStand}{%
  \draw[eq] (0,0) rectangle (2.3,0.2);
  \draw[eq] (0.35,0.2) -- (0.35,4.0) (0.5,0.2) -- (0.5,4.0);
  \draw[eq] (0.25,2.75) rectangle (0.75,3.2);
  \draw[eq] (0.75,2.92) -- (1.6,2.92) (0.75,3.03) -- (1.6,3.03);
  \draw[eq] (1.6,2.7) -- (1.6,3.25) (1.6,2.7) -- (1.95,2.7) (1.6,3.25) -- (1.95,3.25);}
\newcommand{\eqWatchGlass}{\draw[eq] (0,0.35) .. controls (0.5,-0.05) and (1.7,-0.05) .. (2.2,0.35);}
\newcommand{\eqTongs}{%
  \draw[eq] (0,0.1) -- (2.0,0.45) .. controls (2.3,0.5) and (2.45,0.35) .. (2.6,0.15)
            (0,0.0) -- (2.0,-0.35) .. controls (2.3,-0.4) and (2.45,-0.25) .. (2.6,-0.05);}

% A label line from a point on a drawing to text at the side: \eqlab{(x,y)}{(tx,ty)}{text}
\newcommand{\eqlab}[3]{\draw[eqlab] #1 -- #2 node[eqtext, anchor=west] {#3};}
\newcommand{\eqlabl}[3]{\draw[eqlab] #1 -- #2 node[eqtext, anchor=east] {#3};}

% ── Reading scales ──────────────────────────────────────────────────────────
% \readcyl{lowest}{highest}{labelled step}{small divisions per label}{reading}
% A close-up of a measuring cylinder's scale, the water drawn to the reading:
% the bottom of the meniscus sits exactly on the reading.
\pgfkeys{/pgf/number format/.cd, fixed, precision=1}
\newcommand{\readcyl}[5]{%
\begin{tikzpicture}[baseline=(current bounding box.south)]
  \pgfmathsetmacro\rcH{5.0}
  \pgfmathsetmacro\rcY{(#5-#1)*\rcH/(#2-#1)}
  \pgfmathtruncatemacro\rcN{round((#2-#1)/#3*#4)}
  \fill[black!10] (0,-0.3) -- (0,\rcY+0.16) .. controls (0.45,\rcY-0.05) and (1.15,\rcY-0.05) .. (1.6,\rcY+0.16) -- (1.6,-0.3) -- cycle;
  \draw[Ink, line width=0.9pt] (0,\rcY+0.16) .. controls (0.45,\rcY-0.05) and (1.15,\rcY-0.05) .. (1.6,\rcY+0.16);
  \foreach \k in {0,...,\rcN}{
    \pgfmathsetmacro\yy{\k*\rcH/\rcN}
    \pgfmathtruncatemacro\ismaj{mod(\k,#4)==0 ? 1 : 0}
    \ifnum\ismaj=1
      \draw[Ink, line width=0.75pt] (0,\yy) -- (0.6,\yy);
      \pgfmathsetmacro\lv{#1+\k/#4*#3}
      \node[anchor=west, font=\scriptsize\bfseries, inner sep=1pt, fill=white, fill opacity=0.8, text opacity=1] at (0.66,\yy) {\pgfmathprintnumber{\lv}};
    \else
      \draw[Ink, line width=0.4pt] (0,\yy) -- (0.32,\yy);
    \fi}
  \draw[eq] (0,-0.3) -- (0,\rcH+0.35) (1.6,-0.3) -- (1.6,\rcH+0.35);
  \node[font=\scriptsize, anchor=north] at (0.8,-0.35) {scale in mL};
\end{tikzpicture}}
% \readtherm{lowest}{highest}{labelled step}{small divisions per label}{reading}
\newcommand{\readtherm}[5]{%
\begin{tikzpicture}[baseline=(current bounding box.south)]
  \pgfmathsetmacro\rtH{5.0}
  \pgfmathsetmacro\rtY{(#5-#1)*\rtH/(#2-#1)}
  \pgfmathtruncatemacro\rtN{round((#2-#1)/#3*#4)}
  \draw[eq, fill=white] (0,-0.45) circle (0.3);
  \fill[black!75] (0,-0.45) circle (0.22);
  \fill[black!75] (-0.07,-0.3) rectangle (0.07,\rtY);
  \draw[eq] (-0.16,-0.2) -- (-0.16,\rtH+0.3) arc[start angle=180,end angle=0,radius=0.16] -- (0.16,-0.2);
  \foreach \k in {0,...,\rtN}{
    \pgfmathsetmacro\yy{\k*\rtH/\rtN}
    \pgfmathtruncatemacro\ismaj{mod(\k,#4)==0 ? 1 : 0}
    \ifnum\ismaj=1
      \draw[Ink, line width=0.75pt] (0.25,\yy) -- (0.75,\yy);
      \pgfmathsetmacro\lv{#1+\k/#4*#3}
      \node[anchor=west, font=\scriptsize\bfseries, inner sep=1pt] at (0.8,\yy) {\pgfmathprintnumber{\lv}};
    \else
      \draw[Ink, line width=0.4pt] (0.25,\yy) -- (0.52,\yy);
    \fi}
  \node[font=\scriptsize, anchor=north] at (0.3,-0.8) {scale in \textdegree C};
\end{tikzpicture}}
% \balancefig{reading}: an electronic balance showing a reading
\newcommand{\balancefig}[1]{%
\begin{tikzpicture}[baseline=(current bounding box.south)]
  \draw[eq, rounded corners=3pt] (0,0) rectangle (4.2,1.35);
  \draw[eq] (1.9,1.35) -- (1.9,1.55) (2.3,1.35) -- (2.3,1.55);
  \draw[eq, rounded corners=1pt] (0.5,1.55) rectangle (3.7,1.72);
  \draw[eq, fill=LightGrey] (0.35,0.35) rectangle (2.75,1.02);
  \node[font=\headingfont\bfseries\large, anchor=east] at (2.65,0.68) {#1};
  \draw[eq] (3.3,0.68) circle (0.2);
  \node[font=\tiny\bfseries] at (3.3,0.3) {TARE};
\end{tikzpicture}}
% Bunsen burner with a flame: \bunsenflame{safety|heating}
\newcommand{\bunsenflame}[1]{%
\begin{tikzpicture}[baseline=(current bounding box.south)]
  \def\fl{#1}\def\safe{safety}%
  \ifx\fl\safe
    \fill[black!14] (0.8,2.3) .. controls (0.05,3.0) and (0.45,4.3) .. (0.85,5.0) .. controls (1.05,4.4) and (1.55,3.1) .. (0.8,2.3);
    \draw[Ink, line width=0.6pt] (0.8,2.3) .. controls (0.05,3.0) and (0.45,4.3) .. (0.85,5.0) .. controls (1.05,4.4) and (1.55,3.1) .. (0.8,2.3);
  \else
    \fill[white] (0.8,2.3) .. controls (0.45,2.8) and (0.6,3.5) .. (0.8,3.9) .. controls (1.0,3.5) and (1.15,2.8) .. (0.8,2.3);
    \draw[Ink, line width=0.6pt] (0.8,2.3) .. controls (0.45,2.8) and (0.6,3.5) .. (0.8,3.9) .. controls (1.0,3.5) and (1.15,2.8) .. (0.8,2.3);
    \fill[black!30] (0.8,2.3) .. controls (0.65,2.55) and (0.7,2.85) .. (0.8,3.0) .. controls (0.9,2.85) and (0.95,2.55) .. (0.8,2.3);
  \fi
  \draw[eq] (0,0) rectangle (1.6,0.22);
  \draw[eq] (0.58,0.85) -- (0.58,2.3) (1.02,0.85) -- (1.02,2.3);
  \draw[eq] (0.5,0.22) rectangle (1.1,0.85);
  \ifx\fl\safe\draw[Ink, line width=0.5pt, densely dotted] (0.72,0.42) rectangle (0.88,0.62);\else\fill[Ink] (0.72,0.42) rectangle (0.88,0.62);\fi
  \draw[eq] (-0.65,0.06) rectangle (0,0.16);
\end{tikzpicture}}
